% ============================================================================
% SOLUCIONES DEL MÓDULO 7: MACHINE LEARNING PARA BUSINESS INTELLIGENCE
% ============================================================================

\begin{solucion}{m7-1}
\textbf{Parte conceptual.}
\begin{enumerate}[label=(\alph*)]
  \item Estimar el ticket promedio de un cliente nuevo: \textbf{regresión}
        (la respuesta es un número). Regresión lineal o random forest.
  \item Detectar si un correo es una queja: \textbf{clasificación} (queja /
        no queja). Regresión logística o random forest sobre variables del
        texto (palabras clave).
  \item Agrupar sucursales por su mezcla de ventas: \textbf{clustering}.
        k-means con las variables escaladas.
  \item Demanda semanal de pan de muerto: \textbf{serie de tiempo} con
        estacionalidad anual muy marcada. Ingenuo estacional como línea base,
        ETS o ARIMA.
  \item Qué recomendar a quien compra una cafetera: \textbf{reglas de
        asociación} (análisis de canasta) ordenadas por lift.
\end{enumerate}

\textbf{Parte práctica.} Ajustamos el modelo y extraemos el R² ajustado:

\begin{rcode}
library(tidyverse)
ventas <- read_csv("datasets/ventas_retail.csv", show_col_types = FALSE)

modelo <- lm(total ~ descuento_pct + dia_semana, data = ventas)
round(summary(modelo)$r.squared, 4)
round(summary(modelo)$adj.r.squared, 4)
\end{rcode}

\begin{rsalida}
[1] 0.004
[1] -0.0155
\end{rsalida}

El R² es de apenas 0.004 y el R² ajustado es negativo ($-0.0155$): el
descuento y el día de la semana no explican el total de la venta. No hay
señal, porque en el generador de datos el descuento y el día se eligieron al
azar, independientemente del monto.
\end{solucion}

\begin{solucion}{m7-2}
Escribimos cada fórmula directamente con vectores:

\begin{rcode}
real       <- c(120, 95, 150, 80, 110, 130)
pronostico <- c(110, 100, 140, 90, 115, 120)
error <- real - pronostico

mae  <- mean(abs(error))
rmse <- sqrt(mean(error^2))
mape <- mean(abs(error) / real) * 100
r2   <- 1 - sum(error^2) / sum((real - mean(real))^2)

round(c(MAE = mae, RMSE = rmse, MAPE = mape, R2 = r2), 3)
\end{rcode}

\begin{rsalida}
  MAE  RMSE  MAPE    R2 
8.333 8.660 7.500 0.856 
\end{rsalida}

En promedio, el gerente se equivoca por 8.3 mil pesos por sucursal (MAE),
que equivale a 7.5\% del valor real (MAPE). El RMSE (8.66) es apenas mayor que
el MAE porque ningún error es mucho más grande que los demás (todos son de 5 o
10 mil pesos). El R² de 0.856 indica que el pronóstico captura el 85.6\% de las
diferencias entre sucursales.
\end{solucion}

\begin{solucion}{m7-3}
Los datos: 50 fraudes reales, de los que se detectaron 40 (VP = 40, FN = 10);
950 legítimas, de las que 60 se marcaron como fraude (FP = 60, VN = 890).

\begin{rcode}
matriz <- matrix(c(890, 60, 10, 40), nrow = 2,
                 dimnames = list(Prediccion = c("Legítima", "Fraude"),
                                 Real = c("Legítima", "Fraude")))
matriz

vp <- 40; fn <- 10; fp <- 60; vn <- 890
precision    <- vp / (vp + fp)
sensibilidad <- vp / (vp + fn)
round(c(exactitud     = (vp + vn) / (vp + vn + fp + fn),
        precision     = precision,
        sensibilidad  = sensibilidad,
        especificidad = vn / (vn + fp),
        F1 = 2 * precision * sensibilidad / (precision + sensibilidad)), 3)

costo_errores <- fn * 5000 + fp * 100
costo_errores
\end{rcode}

\begin{rsalida}
          Real
Prediccion Legítima Fraude
  Legítima      890     10
  Fraude         60     40
    exactitud     precision  sensibilidad especificidad            F1 
        0.930         0.400         0.800         0.937         0.533 
[1] 56000
\end{rsalida}

La exactitud (93\%) parece alta, pero la precisión es de solo 40\%: de cada 10
alertas, 6 son transacciones legítimas. La sensibilidad de 80\% indica que se
detecta a 4 de cada 5 fraudes. El costo de los errores es de \$56{,}000:
\$50{,}000 por los 10 fraudes no detectados y \$6{,}000 por las 60 revisiones
innecesarias. Con estos costos, conviene bajar el umbral aunque aumenten las
falsas alarmas: cada fraude detectado adicional ahorra lo mismo que 50
revisiones.
\end{solucion}

\begin{solucion}{m7-4}
Calculamos el RFM, las calificaciones con \codigo{ntile()} (recordando que en
recencia menos es mejor) y los niveles con \codigo{case\_when()}:

\begin{rcode}
library(tidyverse)
transacciones <- read_csv("datasets/transacciones.csv",
                          show_col_types = FALSE)
fecha_corte <- as.Date("2024-01-01")

niveles <- transacciones %>%
  group_by(cliente_id) %>%
  summarise(recencia   = as.numeric(fecha_corte - max(fecha)),
            frecuencia = n(),
            monto      = sum(total_transaccion), .groups = "drop") %>%
  mutate(r = ntile(desc(recencia), 5),
         f = ntile(frecuencia, 5),
         m = ntile(monto, 5),
         puntaje = r + f + m,
         nivel = case_when(puntaje >= 13 ~ "Oro",
                           puntaje >= 9  ~ "Plata",
                           TRUE          ~ "Bronce"),
         nivel = factor(nivel, levels = c("Oro", "Plata", "Bronce")))

niveles %>%
  group_by(nivel) %>%
  summarise(clientes = n(),
            monto_prom = round(mean(monto)),
            ingresos = sum(monto), .groups = "drop") %>%
  mutate(pct_ingresos = round(100 * ingresos / sum(ingresos), 1)) %>%
  select(-ingresos)
\end{rcode}

\begin{rsalida}
# A tibble: 3 x 4
  nivel  clientes monto_prom pct_ingresos
  <fct>     <int>      <dbl>        <dbl>
1 Oro          39       8439         34.3
2 Plata        70       5444         39.7
3 Bronce       89       2807         26  
\end{rsalida}

Los 39 clientes Oro (20\% de los 198 compradores) generan el 34.3\% de los
ingresos, con un monto promedio tres veces mayor que el de los Bronce (\$8{,}439
contra \$2{,}807). Los Plata son el grupo con mayor peso en ingresos
(39.7\%): son el candidato natural para programas que los conviertan en Oro.
\end{solucion}

\begin{solucion}{m7-5}
Construimos las variables por cliente, escalamos y calculamos la silueta
promedio para $k$ de 2 a 6:

\begin{rcode}
library(tidyverse)
library(cluster)
transacciones <- read_csv("datasets/transacciones.csv",
                          show_col_types = FALSE)

perfil_cliente <- transacciones %>%
  group_by(cliente_id) %>%
  summarise(monto_total     = sum(total_transaccion),
            ticket_promedio = mean(total_transaccion),
            n_tiendas       = n_distinct(tienda_id),
            n_productos     = n_distinct(producto_id), .groups = "drop")

datos_esc <- scale(perfil_cliente %>% select(-cliente_id))

siluetas <- map_dbl(2:6, function(k) {
  set.seed(123)
  km <- kmeans(datos_esc, centers = k, nstart = 25)
  mean(silhouette(km$cluster, dist(datos_esc))[, "sil_width"])
})
tibble(k = 2:6, silueta = round(siluetas, 3))
k_elegido <- (2:6)[which.max(siluetas)]
k_elegido
\end{rcode}

\begin{rsalida}
# A tibble: 5 x 2
      k silueta
  <int>   <dbl>
1     2   0.322
2     3   0.325
3     4   0.356
4     5   0.326
5     6   0.331
[1] 4
\end{rsalida}

La silueta más alta es para $k = 4$ (0.356), aunque todas están entre 0.32 y
0.36: la estructura de grupos es débil (datos aleatorios). El código guarda el
$k$ elegido en \codigo{k\_elegido} para no escribirlo a mano.

\begin{rcode}
set.seed(123)
km <- kmeans(datos_esc, centers = k_elegido, nstart = 25)

perfil <- perfil_cliente %>%
  mutate(cluster = km$cluster) %>%
  group_by(cluster) %>%
  summarise(clientes = n(),
            monto_total = round(mean(monto_total)),
            ticket = round(mean(ticket_promedio)),
            tiendas = round(mean(n_tiendas), 1),
            productos = round(mean(n_productos), 1), .groups = "drop")

# Nombres según el perfil (reglas de negocio, nunca por número de grupo)
extremos <- range(perfil$monto_total)
perfil %>%
  mutate(segmento = case_when(
    monto_total == extremos[2] ~ "Alto valor",
    monto_total == extremos[1] ~ "Valor mínimo",
    ticket == max(ticket[!monto_total %in% extremos]) ~
      "Ticket alto ocasional",
    TRUE ~ "Valor medio")) %>%
  arrange(desc(monto_total)) %>%
  select(segmento, everything(), -cluster)
\end{rcode}

\begin{rsalida}
# A tibble: 4 x 6
  segmento              clientes monto_total ticket tiendas productos
  <chr>                    <int>       <dbl>  <dbl>   <dbl>     <dbl>
1 Alto valor                  22       10159   1123     6.1       8.5
2 Valor medio                 96        4999    873     4.7       5.6
3 Ticket alto ocasional       39        4756   1547     2.7       3  
4 Valor mínimo                41        1734    597     2.6       2.8
\end{rsalida}

Los nombres salen de reglas aplicadas al perfil: el mayor monto es ``Alto
valor'' (22 clientes que compran en 6 tiendas distintas en promedio), el menor
es ``Valor mínimo'', y entre los dos restantes, el de mayor ticket es ``Ticket
alto ocasional'' (39 clientes que compran pocas veces, en pocas tiendas, pero
con el ticket más alto: \$1{,}547). Este último grupo no se ve en un RFM
clásico y sugiere una estrategia propia: aumentar su frecuencia con
recordatorios, porque cada visita vale mucho.
\end{solucion}

\begin{solucion}{m7-6}
Reproducimos la simulación y la división del texto, y comparamos tres
modelos:

\begin{rcode}
library(tidyverse)
library(caret)

set.seed(123)
n_clientes <- 2000
clientes_churn <- tibble(
  cliente_id       = sprintf("C%04d", 1:n_clientes),
  antiguedad_meses = sample(1:72, n_clientes, replace = TRUE),
  plan             = sample(c("Básico", "Estándar", "Premium"),
                            n_clientes, replace = TRUE,
                            prob = c(0.5, 0.3, 0.2)),
  uso_mensual      = rpois(n_clientes, 8),
  quejas_6m        = rpois(n_clientes, 0.8),
  pagos_tardios    = rpois(n_clientes, 0.4),
  genero           = sample(c("F", "M"), n_clientes, replace = TRUE),
  edad             = sample(18:70, n_clientes, replace = TRUE)
) %>%
  mutate(
    efecto_plan = case_when(plan == "Básico"   ~ 0,
                            plan == "Estándar" ~ -0.6,
                            TRUE               ~ -1.3),
    riesgo = 2.0 - 0.05 * antiguedad_meses +
      1.0 * (antiguedad_meses <= 6) + 1.0 * quejas_6m -
      0.35 * uso_mensual + 0.8 * pagos_tardios + efecto_plan,
    churn = if_else(runif(n_clientes) < plogis(riesgo), "Si", "No"),
    churn = factor(churn, levels = c("No", "Si")),
    plan  = factor(plan, levels = c("Básico", "Estándar", "Premium")),
    cliente_nuevo = as.integer(antiguedad_meses <= 6)   # variable nueva
  ) %>%
  select(-efecto_plan, -riesgo)

set.seed(123)
filas_ent <- createDataPartition(clientes_churn$churn, p = 0.7,
                                 list = FALSE)
ent    <- clientes_churn[filas_ent, ]
prueba <- clientes_churn[-filas_ent, ]

calcular_auc <- function(prob, real, positivo = "Si") {
  p_pos <- prob[real == positivo]
  p_neg <- prob[real != positivo]
  mean(outer(p_pos, p_neg, ">") + 0.5 * outer(p_pos, p_neg, "=="))
}

f_original <- churn ~ antiguedad_meses + plan + uso_mensual + quejas_6m +
  pagos_tardios + genero + edad
f_nuevo    <- update(f_original, . ~ . + cliente_nuevo)
f_justo    <- update(f_nuevo, . ~ . - genero - edad)

modelos <- list(Original = f_original, `Con cliente_nuevo` = f_nuevo,
                `Sin genero ni edad` = f_justo) %>%
  map(~ glm(.x, data = ent, family = binomial))

map_dbl(modelos, ~ calcular_auc(predict(.x, prueba, type = "response"),
                                prueba$churn)) %>%
  round(3)

round(exp(coef(modelos[["Con cliente_nuevo"]])["cliente_nuevo"]), 2)
\end{rcode}

\begin{rsalida}
          Original  Con cliente_nuevo Sin genero ni edad 
             0.860              0.862              0.861 
cliente_nuevo 
         2.36 
\end{rsalida}

Agregar \codigo{cliente\_nuevo} sube el AUC de 0.860 a 0.862 y su razón de
momios es 2.36: estar en los primeros 6 meses multiplica por 2.36 los momios
de abandono, además del efecto de la antigüedad. La mejora en AUC es pequeña
porque la antigüedad lineal ya capturaba parte del efecto, pero la variable
tiene sentido de negocio (reforzar el \emph{onboarding}). Quitar género y edad
casi no cambia el AUC (0.861): el modelo más justo y más simple predice igual
de bien, así que es el que conviene usar.
\end{solucion}

\begin{solucion}{m7-7}
Reproducimos datos, división y modelo logístico, y recalculamos los costos con
el nuevo precio del descuento:

\begin{rcode}
library(tidyverse)
library(caret)

set.seed(123)
n_clientes <- 2000
clientes_churn <- tibble(
  cliente_id       = sprintf("C%04d", 1:n_clientes),
  antiguedad_meses = sample(1:72, n_clientes, replace = TRUE),
  plan             = sample(c("Básico", "Estándar", "Premium"),
                            n_clientes, replace = TRUE,
                            prob = c(0.5, 0.3, 0.2)),
  uso_mensual      = rpois(n_clientes, 8),
  quejas_6m        = rpois(n_clientes, 0.8),
  pagos_tardios    = rpois(n_clientes, 0.4),
  genero           = sample(c("F", "M"), n_clientes, replace = TRUE),
  edad             = sample(18:70, n_clientes, replace = TRUE)
) %>%
  mutate(
    efecto_plan = case_when(plan == "Básico"   ~ 0,
                            plan == "Estándar" ~ -0.6,
                            TRUE               ~ -1.3),
    riesgo = 2.0 - 0.05 * antiguedad_meses +
      1.0 * (antiguedad_meses <= 6) + 1.0 * quejas_6m -
      0.35 * uso_mensual + 0.8 * pagos_tardios + efecto_plan,
    churn = if_else(runif(n_clientes) < plogis(riesgo), "Si", "No"),
    churn = factor(churn, levels = c("No", "Si")),
    plan  = factor(plan, levels = c("Básico", "Estándar", "Premium"))
  ) %>%
  select(-efecto_plan, -riesgo)

set.seed(123)
filas_ent <- createDataPartition(clientes_churn$churn, p = 0.7,
                                 list = FALSE)
ent    <- clientes_churn[filas_ent, ]
prueba <- clientes_churn[-filas_ent, ]

modelo_log <- glm(churn ~ antiguedad_meses + plan + uso_mensual +
                    quejas_6m + pagos_tardios + genero + edad,
                  data = ent, family = binomial)
prob <- predict(modelo_log, prueba, type = "response")
real <- prueba$churn == "Si"

costo_descuento <- 800
costo_vp <- costo_descuento + 0.5 * 3000
costo_fp <- costo_descuento
costo_fn <- 3000

costos <- map_dfr(seq(0.05, 0.95, by = 0.05), function(u) {
  pred <- prob >= u
  tibble(umbral = u, llamadas = sum(pred),
         costo_total = sum(pred & real) * costo_vp +
           sum(pred & !real) * costo_fp + sum(!pred & real) * costo_fn)
})
costos %>% arrange(costo_total) %>% head(3)
\end{rcode}

\begin{rsalida}
# A tibble: 3 x 3
  umbral llamadas costo_total
   <dbl>    <int>       <dbl>
1   0.6        74      381700
2   0.65       61      381800
3   0.7        52      382100
\end{rsalida}

Con un descuento de \$800, el umbral óptimo sube de 0.15 a 0.6 y solo se
llamaría a 74 clientes de prueba (costo total de \$381{,}700). La razón: cada
llamada ahora es más cara y el ahorro al retener a un cliente (la mitad de
\$3{,}000 menos \$800) es mucho menor, así que solo vale la pena llamar a
quien tiene una probabilidad alta de irse. Observa que los tres mejores
umbrales tienen costos casi idénticos: la curva de costos es plana en esa
zona.
\end{solucion}

\begin{solucion}{m7-8}
Reproducimos la serie mensual, la agregamos a trimestres con
\codigo{aggregate()} y comparamos los modelos:

\begin{rcode}
library(tidyverse)
library(forecast)

set.seed(123)
meses <- seq(as.Date("2021-01-01"), as.Date("2025-12-01"), by = "month")
n_meses <- length(meses)
factor_mes <- c(0.88, 0.85, 0.95, 0.97, 1.04, 0.96,
                0.97, 1.00, 0.93, 0.98, 1.15, 1.32)
factor_mes <- factor_mes / mean(factor_mes)
unidades <- round((1500 + 15 * (1:n_meses)) * factor_mes[month(meses)] *
                    (1 + rnorm(n_meses, 0, 0.03)))
serie <- ts(unidades, start = c(2021, 1), frequency = 12)

# Suma por trimestre: frecuencia 4
serie_trim <- aggregate(serie, nfrequency = 4, FUN = sum)
serie_trim

ent    <- window(serie_trim, end = c(2024, 4))
prueba <- window(serie_trim, start = c(2025, 1))

pronosticos <- list(
  `Ingenuo estacional` = snaive(ent, h = 4),
  ETS   = forecast(ets(ent), h = 4),
  ARIMA = forecast(auto.arima(ent), h = 4)
)
comparacion <- map_dfr(pronosticos, function(p) {
  acc <- accuracy(p, prueba)["Test set", ]
  tibble(MAE = acc["MAE"], RMSE = acc["RMSE"], MAPE = acc["MAPE"])
}, .id = "modelo") %>%
  mutate(across(where(is.numeric), ~ round(.x, 1))) %>%
  arrange(MAPE)
comparacion
\end{rcode}

\begin{rsalida}
     Qtr1 Qtr2 Qtr3 Qtr4
2021 4139 4765 4626 5822
2022 4579 5229 5176 6240
2023 4999 5757 5797 7156
2024 5557 6200 6399 7468
2025 6089 6903 6774 8305
# A tibble: 3 x 4
  modelo               MAE  RMSE  MAPE
  <chr>              <dbl> <dbl> <dbl>
1 ETS                 96.2  105.   1.4
2 ARIMA              163.   196.   2.2
3 Ingenuo estacional 612.   636.   8.6
\end{rsalida}

\codigo{aggregate(serie, nfrequency = 4, FUN = sum)} suma los tres meses de
cada trimestre. En trimestres, ETS es el mejor modelo con un MAPE de 1.4\%,
seguido de ARIMA (2.2\%); el ingenuo estacional se equivoca 8.6\% porque no
considera el crecimiento. Reentrenamos el ganador con toda la serie:

\begin{rcode}
# Reentrenamos el ganador (menor MAPE) con toda la serie trimestral
ganador <- comparacion$modelo[1]
modelo_final <- switch(ganador,
                       ETS   = ets(serie_trim),
                       ARIMA = auto.arima(serie_trim),
                       snaive(serie_trim, h = 4))
ganador
forecast(modelo_final, h = 4)
\end{rcode}

\begin{rsalida}
[1] "ETS"
        Point Forecast    Lo 80    Hi 80    Lo 95    Hi 95
2026 Q1       6504.441 6355.576 6653.305 6276.772 6732.109
2026 Q2       7369.522 7200.758 7538.287 7111.419 7627.625
2026 Q3       7303.435 7136.086 7470.783 7047.497 7559.372
2026 Q4       8855.932 8652.896 9058.969 8545.414 9166.450
\end{rsalida}

El código elige el modelo automáticamente con \codigo{switch()} según el
ganador. Para 2026 se esperan 6{,}504 unidades en el primer trimestre y 8{,}856
en el cuarto (entre 8{,}653 y 9{,}059 con 80\% de confianza). Con datos
trimestrales hay menos observaciones (20), pero los errores relativos son
menores porque al sumar tres meses el ruido mensual se compensa.
\end{solucion}

\begin{solucion}{m7-9}
Reproducimos los tickets, marcamos por ticket si tiene cada producto y
contamos:

\begin{rcode}
library(tidyverse)

set.seed(123)
n_tickets <- 3000
tickets_ancho <- tibble(
  ticket_id = 1:n_tickets,
  Laptop    = rbinom(n_tickets, 1, 0.08),
  Monitor   = rbinom(n_tickets, 1, 0.07),
  Impresora = rbinom(n_tickets, 1, 0.05),
  Audifonos = rbinom(n_tickets, 1, 0.15),
  USB       = rbinom(n_tickets, 1, 0.18),
  Papel     = rbinom(n_tickets, 1, 0.20),
  Teclado   = rbinom(n_tickets, 1, 0.08)
) %>%
  mutate(
    Mouse = if_else(Laptop == 1, rbinom(n(), 1, 0.60),
                    rbinom(n(), 1, 0.08)),
    Funda = if_else(Laptop == 1, rbinom(n(), 1, 0.40),
                    rbinom(n(), 1, 0.02)),
    `Cable HDMI` = if_else(Monitor == 1, rbinom(n(), 1, 0.50),
                           rbinom(n(), 1, 0.04)),
    Tinta = if_else(Impresora == 1, rbinom(n(), 1, 0.70),
                    rbinom(n(), 1, 0.06))
  )
tickets <- tickets_ancho %>%
  pivot_longer(-ticket_id, names_to = "producto", values_to = "compro") %>%
  filter(compro == 1) %>%
  select(-compro)

# Una fila por ticket con TRUE/FALSE para los productos de interés
por_ticket <- tickets %>%
  group_by(ticket_id) %>%
  summarise(laptop = "Laptop" %in% producto,
            mouse  = "Mouse" %in% producto,
            funda  = "Funda" %in% producto, .groups = "drop")
n_total <- nrow(por_ticket)

soporte_funda <- mean(por_ticket$funda)
tres <- por_ticket %>%
  reframe(regla = c("Laptop+Mouse -> Funda", "Laptop -> Funda"),
            tickets_x  = c(sum(laptop & mouse), sum(laptop)),
            tickets_xy = c(sum(laptop & mouse & funda),
                           sum(laptop & funda)))
tres %>%
  mutate(soporte   = round(tickets_xy / n_total, 3),
         confianza = round(tickets_xy / tickets_x, 3),
         lift      = round(tickets_xy / tickets_x / soporte_funda, 2))
\end{rcode}

\begin{rsalida}
# A tibble: 2 x 6
  regla                 tickets_x tickets_xy soporte confianza  lift
  <chr>                     <int>      <int>   <dbl>     <dbl> <dbl>
1 Laptop+Mouse -> Funda       136         51   0.026     0.375  5.14
2 Laptop -> Funda             225         82   0.042     0.364  5   
\end{rsalida}

Saber que el cliente ya lleva mouse casi no cambia la probabilidad de que
lleve funda: la confianza pasa de 36.4\% a 37.5\% y el lift de 5.0 a 5.14.
Tiene sentido, porque en la simulación la funda depende solo de la laptop.
Una regla de tres productos solo aporta si su confianza es claramente mayor
que la de la regla más simple; si no, quédate con la simple, que además tiene
más soporte (4.2\% contra 2.6\% de los tickets). Nota: usamos
\codigo{reframe()} en lugar de \codigo{summarise()} porque el resultado tiene
dos filas.
\end{solucion}

\begin{solucion}{m7-10}
Organizamos el reto en funciones: una para simular clientes (la usamos para
el histórico y para el lote nuevo), el entrenamiento con \paquete{caret}, la
elección del modelo y la función de puntuación.

\begin{rcode}
library(tidyverse)
library(caret)
library(randomForest)

simular_clientes <- function(n, prefijo = "C") {
  tibble(
    cliente_id       = sprintf("%s%05d", prefijo, 1:n),
    antiguedad_meses = sample(1:72, n, replace = TRUE),
    plan             = sample(c("Básico", "Estándar", "Premium"), n,
                              replace = TRUE, prob = c(0.5, 0.3, 0.2)),
    uso_mensual      = rpois(n, 8),
    quejas_6m        = rpois(n, 0.8),
    pagos_tardios    = rpois(n, 0.4),
    genero           = sample(c("F", "M"), n, replace = TRUE),
    edad             = sample(18:70, n, replace = TRUE)
  ) %>%
    mutate(
      efecto_plan = case_when(plan == "Básico"   ~ 0,
                              plan == "Estándar" ~ -0.6,
                              TRUE               ~ -1.3),
      riesgo = 2.0 - 0.05 * antiguedad_meses +
        1.0 * (antiguedad_meses <= 6) + 1.0 * quejas_6m -
        0.35 * uso_mensual + 0.8 * pagos_tardios + efecto_plan,
      churn = if_else(runif(n) < plogis(riesgo), "Si", "No"),
      churn = factor(churn, levels = c("Si", "No")),  # "Si" primero: caret
      plan  = factor(plan, levels = c("Básico", "Estándar", "Premium")),
      cuota_mensual = case_when(plan == "Básico"   ~ 199,
                                plan == "Estándar" ~ 349,
                                TRUE               ~ 599)
    ) %>%
    select(-efecto_plan, -riesgo)
}

set.seed(2024)
historico <- simular_clientes(3000)
set.seed(2024)
filas_ent <- createDataPartition(historico$churn, p = 0.7, list = FALSE)
ent    <- historico[filas_ent, ]
prueba <- historico[-filas_ent, ]

# Sin género ni edad (sección de ética)
formula_churn <- churn ~ antiguedad_meses + plan + uso_mensual +
  quejas_6m + pagos_tardios
control <- trainControl(method = "cv", number = 5, classProbs = TRUE,
                        summaryFunction = twoClassSummary)
set.seed(2024)
m_log <- train(formula_churn, data = ent, method = "glm",
               family = binomial, trControl = control, metric = "ROC")
set.seed(2024)
m_rf <- train(formula_churn, data = ent, method = "rf", ntree = 200,
              trControl = control, metric = "ROC",
              tuneGrid = data.frame(mtry = 2:3))

calcular_auc <- function(prob, real, positivo = "Si") {
  p_pos <- prob[real == positivo]
  p_neg <- prob[real != positivo]
  mean(outer(p_pos, p_neg, ">") + 0.5 * outer(p_pos, p_neg, "=="))
}
auc_prueba <- c(
  Logistica = calcular_auc(predict(m_log, prueba, type = "prob")[, "Si"],
                           prueba$churn),
  RandomForest = calcular_auc(predict(m_rf, prueba, type = "prob")[, "Si"],
                              prueba$churn))
round(auc_prueba, 3)

mejor <- list(Logistica = m_log, RandomForest = m_rf)[[
  names(which.max(auc_prueba))]]
dir.create("resultados", showWarnings = FALSE)
saveRDS(mejor, "resultados/modelo_churn_reto.rds")
\end{rcode}

\begin{rsalida}
   Logistica RandomForest 
       0.873        0.843 
\end{rsalida}

La regresión logística gana en prueba (AUC de 0.873 contra 0.843 del random
forest), así que es el modelo que se guarda en
\archivo{resultados/modelo\_churn\_reto.rds}. El código elige al ganador con
\codigo{which.max()}: si mañana gana el bosque, el resto del sistema sigue
funcionando igual, porque ambos modelos de \paquete{caret} se usan con
\codigo{predict(..., type = "prob")}.

\begin{rcode}
puntuar <- function(nuevos, modelo, top = 100) {
  nuevos %>%
    mutate(prob_churn = predict(modelo, newdata = nuevos,
                                type = "prob")[, "Si"],
           perdida_esperada = prob_churn * cuota_mensual * 12) %>%
    arrange(desc(perdida_esperada)) %>%
    slice_head(n = top) %>%
    mutate(prioridad = row_number()) %>%
    select(prioridad, cliente_id, plan, quejas_6m, prob_churn,
           perdida_esperada)
}

set.seed(99)
lote_nuevo <- simular_clientes(1000, prefijo = "N") %>% select(-churn)

modelo <- readRDS("resultados/modelo_churn_reto.rds")
top100 <- puntuar(lote_nuevo, modelo)
write_csv(top100, "resultados/top100_churn.csv")

top100 %>%
  mutate(across(c(prob_churn, perdida_esperada), ~ round(.x, 2))) %>%
  head(5)
nrow(read_csv("resultados/top100_churn.csv", show_col_types = FALSE))
\end{rcode}

\begin{rsalida}
# A tibble: 5 x 6
  prioridad cliente_id plan    quejas_6m prob_churn perdida_esperada
      <int> <chr>      <fct>       <int>      <dbl>            <dbl>
1         1 N00330     Premium         3       0.94            6765.
2         2 N00043     Premium         3       0.89            6399.
3         3 N00814     Premium         2       0.87            6268.
4         4 N00456     Premium         2       0.75            5368.
5         5 N00904     Premium         1       0.67            4787.
[1] 100
\end{rsalida}

La función carga el modelo guardado (como se haría en producción), calcula
la pérdida esperada de cada cliente del lote y exporta los 100 más urgentes;
releer el CSV confirma que tiene 100 filas. Los primeros lugares son clientes
Premium con varias quejas y probabilidades de 0.67 a 0.94. Para ponerlo en
producción faltaría programar el script (Módulo~\ref{cap:reportes}) y
monitorear la deriva de las variables cada semana.
\end{solucion}
